\begin{theorem}[Successor Inequality Reflection]
\label{thm:successor-inequality-reflection}
Let $(P,S,1)$ be a Peano system. For all $a,b \in P$, if $a \neq b$, then
$S(a)\neq S(b)$.
\hyperref[prf:successor-inequality-reflection]{Go to proof.}
\end{theorem}
\end{theorembox}
\LeanFormalizes{thm:successor-inequality-reflection}{lra-lean}{LRA.VolumeII.PeanoSystems.PeanoSystem}{SuccessorInequalityReflection}{pending}

\begin{remark*}[Standard quantified statement]
For a Peano system $(P,S,1)$,
\[
\forall a,b \in P\;(a\neq b \Longrightarrow S(a)\neq S(b)).
\]
\end{remark*}

\begin{remark*}[Predicate reading]
\[
\forall a,b \in P\;(a\neq b \Longrightarrow S(a)\neq S(b)).
\]
\end{remark*}

\begin{remark*}[Negated quantified statement]
\[
\begin{aligned}
&\text{Let } (P,S,1) \text{ be a Peano system.}\\
&\exists a,b \in P\;(a\neq b \land S(a)=S(b)).
\end{aligned}
\]
\end{remark*}

\begin{remark*}[Failure modes]
\begin{description}
  \item[Exposition.]
  Successor inequality reflection fails exactly when two distinct elements
  have the same successor.
  \item[Common successor.]
  \[
  \neg\forall a,b \in P\;(a\neq b \Longrightarrow S(a)\neq S(b)).
  \]
  \[
  \begin{aligned}
  &\exists a,b \in P\;
  \underbrace{(a\neq b \land S(a)=S(b))}_{\text{distinct elements with a common successor}}.
  \end{aligned}
  \]
\end{description}
\end{remark*}

\begin{remark*}[Contrapositive quantified statement]
\[
\begin{aligned}
&\text{Let } (P,S,1) \text{ be a Peano system.}\\
&\forall a,b \in P\;(S(a)=S(b) \Longrightarrow a=b).
\end{aligned}
\]
\end{remark*}

\begin{remark*}[Interpretation]
The successor operation reflects inequality: distinct inputs remain distinct
after one successor step. This is the contrapositive form of successor
injectivity, and it is the form used when an induction step must preserve a
negative equality statement.
\end{remark*}

\begin{dependencies}
\begin{itemize}
  \item \hyperref[def:peano-system]{Peano System}
  \item \hyperref[ax:peano-successor-injective]{Injectivity of Successor}
\end{itemize}
\end{dependencies}

